% 第 6 讲，源文件 src/lecture06.txt（208 行）；R3 块。
\section*{第 6 讲\quad 正合列与 Tate 上同调}

上一讲我们开始讨论一些简单的同调代数；动机是计算某些有限阿贝尔群的阶（特别是 $K^\times/N(L^\times)$，其中 $L/K$ 是局部域的循环扩张）。回忆下面的定义：

\textbf{定义 6.1.} 序列
\[
\cdots\to X_{n-1}\xrightarrow{d_n}X_n\xrightarrow{d_{n+1}}X_{n+1}\to\cdots
\]
称为正合的，若对每个 $n$ 都有 $\ker(d_{n+1})=\operatorname{im}(d_n)$；我们把 $d_i$ 称为微分。

要解这个方程，通常的做法是证明：若 $d_{n+1}$ 杀死某个元素，则该元素落在 $d_n$ 的像中。我们看到，对短正合列 $0\to M\hookrightarrow E\twoheadrightarrow N\to 0$，有 $M=E/N$ 与 $\#E=\#M\cdot\#N$，所以短正合列是度量阿贝尔群大小的有效手段。我们也看到，对任何这样的短正合列与 $n\ge1$，都有长正合列
\begin{equation}\tag{6.1}
0\to M[n]\to E[n]\to N[n]\xrightarrow{\ \delta\ }M/n\to E/n\to N/n\to 0,
\end{equation}
其中回忆
\[
M[n]:=\{x\in M:nx=0\}=\operatorname{Tor}_1(M,\Z/n)=H_1(M\otimes^{\mathbf{L}}\Z/n),
\]
分别表示 $n$-挠子群与第一同调群，$E$、$N$ 类似。边界映射 $\delta$ 把元素 $x\in N[n]$ 提升为 $\tilde x\in E$，使得 $n\tilde x\in M$（因为 $N$ 中 $nx=0$），然后把 $n\tilde x$ 映到它在 $M/n$ 中的等价类。还须验证以下断言：

\textbf{断言 6.2.} 边界映射 $\delta$ 良定义。

\textit{证明.} 设 $\tilde{\tilde x}$ 是 $x$ 的另一个提升。则 $\tilde x-\tilde{\tilde x}\in M$，因为它在 $N$ 中的像为零；于是 $n(\tilde x-\tilde{\tilde x})\in nM$，故在 $M/nM$ 中 $n\tilde x=n\tilde{\tilde x}$。\hfill$\square$

\textbf{断言 6.3.} (6.1) 中的序列正合。

% [原文疑误：原 PDF 第 28 页该证明首局作 "then x̃ ∈ M, and therefore x = 0 in N"，论证跳步
%  （由 δ(x)=nx̃=0 于 M/n 只能得 nx̃∈nM，即 x̃ 与 M 中某元素相差一个 E[n] 元素）；
%  译文按标准论证补全——从而 x 落在 E[n]→N[n] 的像中]
\textit{证明.} 除边界映射处之外，其余各处的正合性都是显然的。若对某个 $x\in N[n]$ 及其提升 $\tilde x\in E$ 有 $\delta(x)=n\tilde x=0$ 于 $M/n$，则 $n\tilde x\in nM$，故 $\tilde x$ 与 $M$ 中某个元素只相差一个 $E[n]$ 中的元素，从而 $x$ 落在 $E[n]\to N[n]$ 的像中。类似地，若 $x\in M/n$ 的像在 $E/n$ 中为零，则 $\tilde x=ny$，其中 $y\in E$，而 $\tilde x$ 是 $x$ 到 $M$ 的提升。投射到 $N$，由正合性得 $0=n\bar y$，从而 $\bar y\in N[n]$。于是再次由正合性得 $ny\in M$，且 $\delta(\bar y)=ny=x$（作为 $M/n$ 中的类），此即所求。\hfill$\square$

我们有下面这个有用的引理：

\textbf{引理 6.4.} 设
\[
0\to X_0\xrightarrow{d_1}X_1\xrightarrow{d_2}\cdots\xrightarrow{d_{n-1}}X_{n-1}\xrightarrow{d_n}X_n\to 0
\]
正合，且所有 $X_i$ 有限。则
\[
\#X_0\cdot\#X_2\cdots=\#X_1\cdot\#X_3\cdots.
\]

\textit{证明.} 对 $n$ 归纳。$n=1$ 时结论显然，故设它对 $n-1$ 成立。构造正合列
\[
0\to X_0\to\cdots\to X_{n-2}\to\ker(d_{n-1})\to 0
\]
与
\[
0\to\ker(d_{n-1})\to X_{n-1}\xrightarrow{d_{n-1}}X_n\to 0
\]
（后者的正合性来自 $d_{n-1}$ 满射到 $X_n$，因为原序列以 $X_n\to 0$ 结尾）。设 $n$ 为偶数。对第一个正合列用归纳假设，得
% [原文疑误：抽出文本作 "#X0···#Xn = #X0···#X_{n−1}·#X_{n−1}/#Im(d_{n−1}) = #X1···#Im(d_{n−1})·X_{n−1}/#Im(d_{n−1})"，
%  其中 #X_{n−1} 出现两次、指标链自相矛盾（按原文第一行，两端相差一个 #X_{n−1}/#Im 因子）。
%  此处按标准证明重写：改用 ker(d_{n−1}) 而非 Im(d_{n−1})，并把归纳用在长度为 n 的截断序列上。]
\[
\#X_0\#X_2\cdots\#X_{n-2}=\#X_1\#X_3\cdots\#X_{n-3}\cdot\#\ker(d_{n-1}),
\]
又由第二个正合列有 $\#X_{n-1}=\#\ker(d_{n-1})\cdot\#X_n$。于是
\[
\#X_0\#X_2\cdots\#X_n
=\#X_1\#X_3\cdots\#X_{n-3}\cdot\#\ker(d_{n-1})\cdot\#X_n
=\#X_1\#X_3\cdots\#X_{n-1},
\]
即所欲证。$n$ 为奇数的情形类似。\hfill$\square$

\textbf{定义 6.5.} 设 $M$ 为阿贝尔群，且 $M/n$ 与 $M[n]$ 都有限。则
\[
\chi(M):=\chi_n(M):=\frac{\#(M/n)}{\#(M[n])}
\]
称为 $M$ 的欧拉示性数。

\textbf{例 6.6.} (1) 若 $M$ 有限，则 $\chi(M)=1$。为此注意
\[
0\to M[n]\to M\xrightarrow{\ n\ }M\to M/n\to 0
\]
正合，于是由引理 6.4 得 $\#(M[n])\cdot\#M=\#M\cdot\#(M/n)$。

(2) 若 $M=\Z$，则 $\chi(M)=n$，因为 $M[n]=0$，而 $M/n=\Z/n\Z$ 的阶为 $n$。

下面这个引理是关于欧拉示性数的一个重要事实：

\textbf{引理 6.7.} 对短正合列 $0\to M\to E\to N\to 0$，若三个阿贝尔群中有两个的 $\chi$ 存在，则第三个的 $\chi$ 也存在，且 $\chi(M)\cdot\chi(N)=\chi(E)$；这里「存在」指的是（比方说对 $M$）$M/n$ 与 $M[n]$ 都有限。

\textit{证明.} 有正合列
\[
0\to M[n]\to E[n]\to N[n]\to M/n\to E/n\to N/n\to 0.
\]
更一般地，注意若 $X_{n-1}\xrightarrow{d_{n-1}}X_n\xrightarrow{d_n}X_{n+1}$ 正合，而 $X_{n-1}$ 与 $X_{n+1}$ 有限，则 $X_n$ 也有限，因为有短正合列
\[
0\to\operatorname{im}(d_{n-1})=\ker(d_n)\to X_n\to\operatorname{im}(d_n)\to 0,
\]
其中两端群有限，故 $\#X_n=\#\ker(d_n)\cdot\#\operatorname{im}(d_n)$ 也有限。于是序列中所有群都有限，且由引理 6.4，
\[
\#(M[n])\cdot\#(N[n])\cdot\#(E/n)=\#(E[n])\cdot\#(M/n)\cdot\#(N/n),
\]
由此即得所要的表达式。\hfill$\square$

作为应用，我们来计算 $\#(K^\times/(K^\times)^n)$。注意
\[
\chi(K^\times)=\frac{\#(K^\times/(K^\times)^n)}{\#(K^\times[n])},
\]
其中分母是 $K$ 中 $n$ 次单位根的个数。此外我们还有正合列
\[
0\to\mathcal{O}_K^\times\to K^\times\xrightarrow{\ v\ }\Z\to 0,
\]
于是由引理 6.7，$\chi(K^\times)=\chi(\mathcal{O}_K^\times)\chi(\Z)=n\chi(\mathcal{O}_K^\times)$。因此我们真正想算的是 $\chi(\mathcal{O}_K^\times)$。

一个好用的启发是：若 $\mathcal{O}_K^\times$ 包含某个开子群（即有限指数子群）$\Gamma$，则 $\Gamma\cong\mathcal{O}_K^+$；当 $\operatorname{char}(K)=0$ 时这可由 $p$ 进指数映射得到。于是
\begin{equation}\tag{6.2}
\chi(\mathcal{O}_K^\times)=\chi(\Gamma)\chi(\mathcal{O}_K^\times/\Gamma)=\chi(\Gamma)=\chi(\mathcal{O}_K)
\end{equation}
（末一个 $\chi$ 是对加法群取的），因为按假设 $\mathcal{O}_K^\times/\Gamma$ 有限。又加法群有 $\mathcal{O}_K[n]=0$（因为 $\mathcal{O}_K$ 是整环），且 $\chi(\mathcal{O}_K)=\#(\mathcal{O}_K/n)=|n|_K^{-1}$，其中 $|x|_K:=q^{-v(x)}$ 表示 $K$ 中正规化的绝对值（即对素元 $\pi$ 有 $v(\pi)=1$），$q$ 表示剩余域的阶。所得的公式
\begin{equation}\tag{6.3}
\#(K^\times/(K^\times)^n)=\frac{n\cdot\#(K^\times[n])}{|n|_K}
\end{equation}
重现了习题集 1 第 1(b) 题中当 $n=2$ 时已证的结果（不过同样的方法对一般的 $n$ 也适用）。不用指数映射的证明用到了这样的事实：对充分大的 $N$，
\[
1+\mathfrak{p}^N\xrightarrow{\ x\mapsto x^n\ }1+\mathfrak{p}^{N+v(n)}
\]
是同构（这可用滤链证明；它是我们先前那个加性陈述的乘性版本）。

现在我们引入循环群的 Tate 上同调。

\textbf{定义 6.8.} 若 $G$ 是（不必有限的）群，则 $G$-模 $A$ 是一个阿贝尔群，$G$ 通过群自同构作用在 $A$ 上。等价地，存在同态 $G\to\operatorname{Aut}(A)$，且该作用满足
\begin{enumerate}[label=(\arabic*),leftmargin=2.2em,itemsep=0pt,topsep=2pt]
\item $g\cdot(a+b)=g\cdot a+g\cdot b$，
\item $(gh)\cdot a=g\cdot(h\cdot a)$，
\end{enumerate}
对一切 $g,h\in G$、$a,b\in A$。

\textbf{例 6.9.} 若 $L/K$ 是域扩张，$G:=\Gal(L/K)$，则 $L$ 与 $L^\times$ 都是 $G$-模，因为域自同构同时保持这两种运算。这将是我们主要关心的例子。

现在假定 $G$ 有限，$A$ 为 $G$-模。

\textbf{定义 6.10.} 第一个 Tate 上同调群为
\[
\widehat{H}^0(G,A):=A^G/N(A),
\]
其中
\[
A^G:=\{a\in A:\ g\cdot a=a\ \text{对一切}\ g\in G\}
\]
表示不动点集。

注意范映射定义为
\[
N\colon A\to A,\qquad a\mapsto\sum_{g\in G}g\cdot a,
\]
所以我们确实需要 $G$ 有限的假设。此外，这个表达式表明范映射经由 $A^G\subseteq A$ 分解，即其像落在 $A^G$ 中。

\textbf{例 6.11.} (1) 回到例 6.9 并取 $A=L$，有 $A^G=K$，而 $N\colon L\to K$ 是域迹，故 $\widehat H^0(L/K)=K/T(L)=0$，因为 $L/K$ 必可分。

(2) 若 $A=L^\times$，则 $(L^\times)^G=K^\times$，且 $\widehat H^0(L^\times)=K^\times/N(L^\times)$。于是我们先前的问题现在重新表述为：对局部域的循环扩张 $L/K$，计算 $\widehat H^0(G,L^\times)$。

(3) 若 $A$ 是任意阿贝尔群，则称 $G$ 平凡地作用在 $A$ 上，若对一切 $g\in G$、$a\in A$ 有 $g\cdot a=a$。此时 $\widehat H^0(G,A)=A/(\#G)A$。于是 Tate 上同调的概念完全推广了我们先前的讨论。

\textbf{定义 6.12.} $G$-模的映射（或称 $G$-态射，或其他任何合理的叫法）$A\xrightarrow{f}B$ 是保持 $G$ 作用的群同态，即对一切 $g\in G$、$a\in A$ 有 $f(g\cdot a)=g\cdot f(a)$。

$G$-模的（短）正合列是指阿贝尔群的（短）正合列，但其中所有映射都是 $G$-态射。

\textbf{例 6.13.} $1\to\mathcal{O}_L^\times\to L^\times\xrightarrow{v}\Z\to1$ 是 $G$-模的短正合列，其中 $G:=\Gal(L/K)$，而 $G$ 平凡地作用在 $\Z$ 上、并通过 Galois 作用作用在 $\mathcal{O}_L^\times$ 上。

现在设 $0\to A\to B\to C\to 0$ 是 $G$-模的短正合列。则我们得到正合列
\begin{equation}\tag{6.4}
\widehat H^0(G,A)\xrightarrow{\ \alpha\ }\widehat H^0(G,B)\xrightarrow{\ \beta\ }\widehat H^0(G,C),
\end{equation}
其中 $\alpha$ 不必是单射（正如上一讲群作用平凡时所见到的那样），$\beta$ 不必是满射。这是因为 Tate 上同调涉及两种操作：其一取不动点，是左正合而非右正合的；其二取商，是右正合而非左正合的。

现在设 $G=\Z/n\Z$，$\sigma\in G$ 为生成元（即 $1$）。

\textbf{定义 6.14.} 第二个 Tate 上同调群为
\[
\widehat H^1(G,A):=\ker(N\colon A\to A)/(1-\sigma)A.
\]

注意我们之所以取商，是因为对任意形如 $x:=y-\sigma y$（$y\in A$）的元素，都有
\[
N(x)=x+\sigma x+\cdots+\sigma^{n-1}x
=y-\sigma y+\sigma y-\sigma^2 y+\cdots+\sigma^{n-1}y-\underbrace{\sigma^{n}y}_{y}=0,
\]
而我们希望把这些平凡情形从核中略去。

现在断言：对正合列 $0\to A\to B\to C\to 0$，存在正合列
\begin{equation}\tag{6.5}
\widehat H^0(A)\to\widehat H^0(B)\to\widehat H^0(C)\xrightarrow{\ \delta\ }
\widehat H^1(A)\to\widehat H^1(B)\to\widehat H^1(C),
\end{equation}
其边界映射 $\delta$ 把任意 $x\in C^G/N(C)$ 提升为 $\tilde x\in B$，然后取 $(1-\sigma)\tilde x$。由于 $x\in C^G$，在 $C$ 中有 $(1-\sigma)x=0$，故 $(1-\sigma)\tilde x\in A$。此外，$(1-\sigma)\tilde x$ 显然被 $A$ 中的范映射杀死，因此它给出 $\widehat H^1(G,A)$ 中的一个类。同样地，我们验证以下两点：

\textbf{断言 6.15.} 边界映射 $\delta$ 良定义，即它不依赖于 $\tilde x$ 的选取。

\textit{证明.} 若 $\tilde{\tilde x}$ 是另一个提升，则 $\tilde x-\tilde{\tilde x}\in A$（因为 $C\cong B/A$），故 $(1-\sigma)(\tilde x-\tilde{\tilde x})$ 在 $\widehat H^1(G,A)$ 中为零。\hfill$\square$

\textbf{断言 6.16.} 序列 (6.5) 可延展为正合的。

\textit{证明.} 与前面一样，我们只在边界映射处验证。取 $x\in B^G/N(B)$，它在 $\widehat H^1(A)$ 中的像是 $(1-\sigma)x=0$。若 $x\in\ker(\delta)$，则对 $x$ 的某个提升 $\tilde x$ 有 $\tilde x\in B^G$，从而落在 $\widehat H^0(B)$ 中。

取 $x\in C^G/N(C)$，它在 $\widehat H^1(A)$ 中的像是 $(1-\sigma)\tilde x$，其中 $\tilde x$ 是 $x$ 到 $B$ 的提升；按定义它在 $\widehat H^1(B)$ 中被杀死。若 $x\in\widehat H^1(A)$ 在 $\widehat H^1(B)$ 中为零，则 $x\in(1-\sigma)B$，从而 $x\in\operatorname{im}(\delta)$。\hfill$\square$
